% !PRINTED 37
\weberpage{37}
Aus (12) folgt ferner
\[
1 = \Theta^a = A^{\frac{\lambda a}{b}};
\]
also muss $\lambda a : b$ ein Vielfaches von $a$ und folglich $\lambda$ ein Vielfaches von $b$ sein. Wenn wir daher
\weq{13}{B = \Theta A^{-\frac{\lambda}{b}}}
setzen, so ist $B^b = 1$, und zugleich ist $B^b$ die niedrigste Potenz von $B$, die einer Potenz von $A$ gleich wird, weil, wenn $B^{b'}$ eine Potenz von $A$ ist, dasselbe nach (13) auch von $\Theta^{b'}$ gilt. Wir nehmen das Element $\Theta$ so gewählt an, dass $b$ so gross als möglich wird. Dann ist auch für jedes andere Element $\Theta_1$ aus $S$ immer $\Theta_1^b$ eine Potenz von $A$, deren Exponent durch $b$ theilbar ist.

Ist nun
\[
\begin{gathered}
\alpha = 0, 1, 2 \ldots a-1 \\
\beta = 0, 1, 2 \ldots b-1,
\end{gathered}
\]
so sind die Elemente
\weq{14}{A^\alpha B^\beta}
in der Anzahl $ab$ alle von einander verschieden. Sie bilden einen in $S$ enthaltenen Theiler $S'$, der durch die Basis $A, B$ dargestellt ist. Ist $S'$ mit $S$ identisch, so sind wir am Ziele.

Ist aber $S$ durch (14) noch nicht erschöpft, so fahren wir in derselben Weise fort.

Wir wollen durch Anwendung der vollständigen Induction gleich allgemein schliessen.

Es sei ein Theiler $S_{\nu-1}$ von $S$ ermittelt, der durch eine Basis in der Form dargestellt ist
\weq{15}{S_{\nu-1} = A_1^{\alpha_1} A_2^{\alpha_2} \ldots A_{\nu-1}^{\alpha_{\nu-1}},}
von dem wir folgende Voraussetzungen machen:

1) Die Grade von $A_1, A_2 \ldots A_{\nu-1}$ seien $a_1, a_2 \ldots a_{\nu-1}$, und es sei
\[
a_1 \geqq a_2 \geqq \ldots \geqq a_{\nu-1}.
\]

2) Für jedes in $S$ enthaltene Element $\Theta$ sei
\weq{16}{\Theta^{a_{\nu-1}} = A_1^{\lambda_1} A_2^{\lambda_2} \ldots A_{\nu-2}^{\lambda_{\nu-2}},}
d. h. in $S_{\nu-2}$ enthalten, und die Exponenten $\lambda_1, \lambda_2 \ldots \lambda_{\nu-2}$ seien durch $a_{\nu-1}$ theilbar.

Die oben abgeleitete Gruppe $S'$ genügt diesen Forderungen, wenn $\nu = 3$ und $a_{\nu-1} = b$ gesetzt wird, und es ist nun zu zeigen,
%CONT>>
